% Appendix: Additional Vision Results
% Contains detailed figures from Paper Figures 4-7 reproduction

\section{Additional Vision Results}
\label{app:vision_additional}

This appendix provides additional visualizations from our reproduction of the paper's vision experiments (Section 4).

\subsection{Noise Sweep Eigenvectors (Paper Figure 4)}
\label{app:noise_sweep}

Figure~\ref{fig:app_noise_eigenvectors} shows how the top eigenvector for digit 5 changes across noise augmentation levels.
Higher noise produces more diffuse, averaged eigenvectors that better capture the canonical digit shape.

\begin{figure}[h]
    \centering
    \includegraphics[width=0.9\linewidth]{figures/vision/figure_4_noise_eigenvectors.pdf}
    \caption{Top eigenvector for digit 5 across noise levels ($\sigma = 0.0$ to $0.5$).
    Accuracy improves with noise despite spreading variance across more eigenvectors.}
    \label{fig:app_noise_eigenvectors}
\end{figure}

\subsection{Challenge Task Regularization Variants (Paper Figure 6)}
\label{app:challenge_variants}

We trained challenge task models under four regularization configurations to analyze how regularization affects the eigendecomposition structure.
Figure~\ref{fig:app_challenge_decay} compares eigenvalue decay across these variants.

\begin{figure}[h]
    \centering
    \includegraphics[width=0.95\linewidth]{figures/vision/figure_6_challenge_eigenvalue_decay_by_reg.pdf}
    \caption{Challenge task eigenvalue decay by regularization configuration.
    Weight decay produces the sharpest spectral decay, concentrating information in fewer eigenvectors.}
    \label{fig:app_challenge_decay}
\end{figure}

\subsection{Eigenvector Quality Comparison}
\label{app:eigenvector_quality}

Figure~\ref{fig:app_eigenvectors_noise} shows the effect of noise-only regularization on eigenvector interpretability.

\begin{figure}[h]
    \centering
    \includegraphics[width=0.95\linewidth]{figures/vision/eigenvectors_noise.pdf}
    \caption{Top eigenvectors with noise augmentation only ($\sigma=0.5$, $\lambda=0.0$).
    Noise alone produces recognizable digit patterns, though less sharp than with weight decay.}
    \label{fig:app_eigenvectors_noise}
\end{figure}

\subsection{Model Size Analysis (Paper Figure 5)}
\label{app:size_analysis}

Our reproduction of Figure 5 examines how eigenvector similarity and truncation accuracy vary with model size.
Figure~\ref{fig:app_truncation} shows classification error under rank truncation across model sizes.

\begin{figure}[h]
    \centering
    \begin{subfigure}{0.48\textwidth}
        \includegraphics[width=\linewidth]{figures/vision/figure_5a_similarity.pdf}
        \caption{Eigenvector similarity across seeds}
        \label{fig:app_similarity}
    \end{subfigure}
    \hfill
    \begin{subfigure}{0.48\textwidth}
        \includegraphics[width=\linewidth]{figures/vision/figure_5b_truncation.pdf}
        \caption{Truncation error by model size}
        \label{fig:app_truncation_error}
    \end{subfigure}
    \caption{Paper Figure 5 reproduction.
    \textbf{(a)} Top eigenvectors are highly consistent across random seeds, especially for larger models.
    \textbf{(b)} Truncation error decreases with model size; ~10 eigenvectors per class suffice for near-full accuracy.}
    \label{fig:app_truncation}
\end{figure}

\subsection{Cross-Dataset Comparison: MNIST vs Fashion-MNIST}
\label{app:cross_dataset}

Figure~\ref{fig:app_fashion} shows eigenvector quality on Fashion-MNIST under noise-only regularization.

\begin{figure}[h]
    \centering
    \includegraphics[width=0.95\linewidth]{figures/vision/fashion_eigenvectors_noise.pdf}
    \caption{Fashion-MNIST top eigenvectors with noise augmentation ($\sigma=0.5$).
    Class labels: T-shirt, Trouser, Pullover, Dress, Coat, Sandal, Shirt, Sneaker, Bag, Boot.}
    \label{fig:app_fashion}
\end{figure}
